%%% FORMATO E IDIOMA %%%
    \documentclass[letterpaper,DIV=15,12pt]{scrartcl}
    \usepackage[spanish,mexico,shorthands=off,es-lcroman]{babel}
    \usepackage{scrlayer-scrpage}
        \clearpairofpagestyles
        \ihead{\footnotesize \textit{Teoría de Conjuntos III}}
        \ohead{\footnotesize \textit{2026-2}}
        \cfoot{\normalfont\thepage}
        \addtokomafont{title}{\bfseries \rmfamily}
        \setlength{\headsep}{5pt}
        \newpairofpagestyles{beginstyle}{
            \clearpairofpagestyles
            \KOMAoptions{headsepline=false}
            \cfoot{\footnotesize \pagemark}
            \cfoot{\normalfont\thepage}
        }
        \addtokomafont{section}{\raggedleft\rmfamily}
        \addtokomafont{subsection}{\raggedleft\rmfamily}
        \addtokomafont{subsubsection}{\raggedleft\rmfamily}
        \setlength{\headsep}{8pt}
        \setlength{\footskip}{20pt}
        \setlength{\textheight}{600pt}
%%% PAQUETES %%%
    \usepackage{ifthen}
    \usepackage[shortlabels]{enumitem}
        \setenumerate[1]{label=\MakeLowercase{\roman*}), ref=\roman*}
        \setenumerate[2]{label=\MakeLowercase{\alph*}), ref=\alph*}
    \usepackage{amsmath}
    \usepackage{amsthm}\usepackage[T1]{fontenc}
    \usepackage{stix2}
	\renewcommand{\qedsymbol}{$\blacksquare$}
	\addto\captionsspanish{\renewcommand{\proofname}{\normalfont\bfseries Demostración}}	
	\newenvironment{solucion}{\renewcommand{\qedsymbol}{$\boxtimes$}\begin{proof}[\normalfont\bfseries Solución]}{\end{proof}}
%%% E J E R C I C I O S %%%
    \newcounter{Ejer}
    \newcounter{Pts}
    \setcounter{Ejer}{1}
    \setcounter{Pts}{1}
    \newcommand{\pts}{}
    \newenvironment{ejercicio}[1]{\noindent
        \ifthenelse{\equal{#1}{1}}{\renewcommand{\pts}{\textbf{(#1 pt)}}}{\renewcommand{\pts}{\textbf{(#1 pts)}}}\textbf{Ej. \theEjer} \pts\stepcounter{Ejer}}{}
%%% C O M A N D O S %%%
    \renewcommand{\emptyset}{\varnothing}
    \newcommand{\tq}{:}
%% D O C U M E N T O %%
\begin{document}
\thispagestyle{plain}

       \begin{center}
            {\fontsize{30}{60}\rmfamily \textbf{Ejercicio semanal 3}} \\ \vspace{.2cm} Teoría de Conjuntos III, 2026-2
       \end{center}
        \begin{flushright}
            \footnotesize \hfill Profesor: Luis Jesús Turcio Cuevas.\\
            \hfill Ayudante: Hugo Víctor García Martínez.
        \end{flushright}
   
       \textit{\textbf{\textsc{Instrucciones.}} Esta tarea es \textbf{individual} y deberá ser entregada \textbf{presencial y personalmente} el día \textbf{lunes 2 de marzo} al inicio la clase.}\vspace{.3cm}

    \begin{ejercicio}{10}
        Sean $\mathscr{I}:=\{ [a,b) \mid a,b \in \mathbb{R} \}$ y $B:=\{ \bigcup F \mid F \subseteq \mathscr{I} \}$.
        \begin{enumerate}
          \item ¿Es $(B,\subseteq)$ un álgebra de Boole?
          \item Prueba que para ningún conjunto $X$ existe una biyección $f:\mathscr{P}(X) \to B$ de modo que para cualesquiera $A,B \subseteq X$ ocurra $A \subseteq B$ si y sólo si $f(A) \subseteq f(B)$. Es decir, muestra que $(\mathscr{P}(X),\subseteq)$ y $(B,\subseteq)$ no son órdenes isomorfos.
        \end{enumerate}
    \end{ejercicio}

    \begin{solucion}\textit{\textbf{(i)}}
    Sean $U,V \in B$, entonces existen $F,G \subseteq \mathscr{I}$ tales que $U=\bigcup F$ y $V=\bigcup G$. Como $F \cup G \subseteq \mathscr{I}$ y $U \cup V = \bigcup (F \cup G)$, entonces $U \cup V \in B$. Ahora, si $[x,y) \in F$ y $[a,b) \in G$, entonces
    \[ [x,y) \cap [a,b) = [\max\{x,a\}, \min\{y,b\}) \in \mathscr{I} \, ; \]
    por lo tanto $H = \{ f \cap g \tq (f,g) \in F \times G \} \subseteq \mathscr{I}$, y así, $U \cap V = \bigcup H \in B$.

    Luego, $B$ tiene ínfimos y supremos dos a dos y son la intersección y la unión. De ser $B$ un álgebra de Boole, sería necesario que para cada $U \in B$, $U U^* = 0_B$ y $U + U^*=1_B$; esto es, $U \cap U^* = \emptyset$ y $U \cup U^* = \mathbb{R}$, de donde $U^* = \mathbb{R} \setminus U$, probando\footnote{Es necesario demostrar que los complementos de $B$ deberían ser, en caso de existir, los complementos respecto $\mathbb{R}$. El hecho de que el orden sea contención, no siempre implica esto. Por ejemplo, los abiertos regulares $\operatorname*{RO}$, ordenados por contención, forman un álgebra de Boole, hay ínfimos, supremos y complementos, pero no todos ellos coinciden con los ``de siempre''.} que $B$ debería ser cerrado bajo complementos respecto $\mathbb{R}$.

    Nótese que $\mathbb{R} \setminus \{0\} = \bigcup \big( \big\{ [x-1,x) \tq x < 0 \big\} \cup \big\{ [x,x+1) \tq x > 0 \big\} \big) \in B$. Sin embargo, $\{ 0 \}$ no es elemento de $B$, pues cada elemento de $B$ debe contener un intervalo no vacío de la forma $[a,b)$, como estos intervalos tienen más de un elemeno, no es posible que $\{0\} \in B$. Por lo tanto, $B$ no es cerrado bajo complementos respecto $\mathbb{R}$.
    \end{solucion}
\end{document}